\bibitem{DelbartBGNPBCBDTJLTU2018}
F. Delbart {\em et al.}, An allosteric binding site of the $\alpha$7 nicotinic acetylcholine receptor revealed in a humanized acetylcholine-binding protein, {\em Journal of Biological Chemistry} {\bf 293} (7), 2534--2545 (2018).

\bibitem{CroxenNBOCVN1997}
R. Croxen {\em et al.}, Mutations in Different Functional Domains of the Human Muscle Acetylcholine Receptor   Subunit in Patients with the Slow-Channel congenital Myasthenic Syndrome, {\em Human Molecular Genetics} {\bf 6} (5), 767--774 (1997).

\bibitem{MiloneWOFPBSE1997}
M. Milone {\em et al.}, Slow-Channel Myasthenic Syndrome Caused By Enhanced Activation, Desensitization, and Agonist Binding Affinity Attributable to Mutation in the M2 Domain of the Acetylcholine Receptor $\alpha$ Subunit, {\em The Journal of Neuroscience} {\bf 17} (15), 5651--5665 (1997).

\bibitem{ShenDSE2006}
X. Shen {\em et al.}, Slow-channel mutation in acetylcholine receptor $\alpha$M4 domain and its efficient knockdown, {\em Annals of Neurology} {\bf 60} (1), 128--136 (2006).

\bibitem{NewlandBVN1995}
C. F. Newland {\em et al.}, Functional and non-functional isoforms of the human muscle acetylcholine receptor., {\em The Journal of Physiology} {\bf 489} (3), 767--778 (1995).

\bibitem{OhnoHMBBSE1995}
K. Ohno {\em et al.}, Congenital myasthenic syndrome caused by prolonged acetylcholine receptor channel openings due to a mutation in the M2 domain of the epsilon subunit., {\em Proceedings of the National Academy of Sciences} {\bf 92} (3), 758--762 (1995).

\bibitem{OhnoQMWHNBPFCSE1997}
K. Ohno {\em et al.}, Congenital Myasthenic Syndromes due to Heteroallelic Nonsense/Missense Mutations in the Acetylcholine Receptor   Subunit Gene: Identification and Functional Characterization of Six New Mutations, {\em Human Molecular Genetics} {\bf 6} (5), 753--766 (1997).

\bibitem{EngelOBSG1996}
A. G. Engel {\em et al.}, End-plate acetylcholine receptor deficiency due to nonsense mutations in the $\varepsilon$ subunit, {\em Annals of Neurology} {\bf 40} (5), 810--817 (1996).

\bibitem{ShenOFTBDPE2002}
X. Shen {\em et al.}, Congenital myasthenic syndrome caused by low-expressor fast-channel AChR $\delta$ subunit mutation, {\em Neurology} {\bf 59} (12), 1881--1888 (2002).

\bibitem{BrownlowWCBNLVNB2001}
S. Brownlow {\em et al.}, Acetylcholine receptor $\delta$ subunit mutations underlie a fast-channel myasthenic syndrome and arthrogryposis multiplex congenita, {\em Journal of Clinical Investigation} {\bf 108} (1), 125--130 (2001).

\bibitem{GomezMGLTCLMK1996}
C. M. Gomez {\em et al.}, A $\beta$-subunit mutation in the acetylcholine receptor channel gate causes severe slow-channel syndrome, {\em Annals of Neurology} {\bf 39} (6), 712--723 (1996).

\bibitem{DurmusSSKPOBDE2018}
H. Durmus {\em et al.}, Congenital myasthenic syndromes in Turkey: Clinical clues and prognosis with long term follow-up, {\em Neuromuscular Disorders} {\bf 28} (4), 315--322 (2018).

\bibitem{BEESONWECBBNSHSCV2003}
D. BEESON {\em et al.}, Structural Abnormalities of the AChR Caused by Mutations Underlying Congenital Myasthenic Syndromes, {\em Annals of the New York Academy of Sciences} {\bf 998} (1), 114--124 (2003).

\bibitem{MasudaSIMEO2008}
A. Masuda {\em et al.}, hnRNP H enhances skipping of a nonfunctional exon P3A in CHRNA1 and a mutation disrupting its binding causes congenital myasthenic syndrome, {\em Human Molecular Genetics} {\bf 17} (24), 4022--4035 (2008).

\bibitem{FaberMVBVDKWBB2009}
C. G. Faber {\em et al.}, AChR deficiency due to $\varepsilon$-subunit mutations: two common mutations in the Netherlands, {\em Journal of Neurology} {\bf 256} (10), 1719--1723 (2009).

\bibitem{KouvatsosGCPT2016}
N. Kouvatsos {\em et al.}, Crystal structure of a human neuronal nAChR extracellular domain in pentameric assembly: Ligand-bound $\alpha$2 homopentamer, {\em Proceedings of the National Academy of Sciences} {\bf 113} (34), 9635--9640 (2016).

\bibitem{AridonMDBDPPMPPCCPBGC2006}
P. Aridon {\em et al.}, Increased Sensitivity of the Neuronal Nicotinic Receptor $\alpha$2 Subunit Causes Familial Epilepsy with Nocturnal Wandering and Ictal Fear, {\em The American Journal of Human Genetics} {\bf 79} (2), 342--350 (2006).

\bibitem{ContiACBMMAMLPRBG2015}
V. Conti {\em et al.}, Nocturnal frontal lobe epilepsy with paroxysmal arousals due to CHRNA2 loss of function, {\em Neurology} {\bf 84} (15), 1520--1528 (2015).

\bibitem{DashL2014}
B. Dash {\em et al.}, Analysis of rare variations reveals roles of amino acid residues in the N-terminal extracellular domain of nicotinic acetylcholine receptor (nAChR) alpha6 subunit in the functional expression of human alpha6*-nAChRs, {\em Molecular Brain} {\bf 7} (1), 35 (2014).

\bibitem{DashBCL2012}
B. Dash {\em et al.}, Modulation of recombinant, $\alpha$2*, $\alpha$3* or $\alpha$4*-nicotinic acetylcholine receptor (nAChR) function by nAChR $\beta$3 subunits*, {\em Journal of Neurochemistry} {\bf 121} (3), 349--361 (2012).

\bibitem{HodaWBS2009}
J. Hoda {\em et al.}, Pleiotropic functional effects of the first epilepsy-associated mutation in the human CHRNA2 gene, {\em FEBS Letters} {\bf 583} (10), 1599--1604 (2009).

\bibitem{DashLL2014}
B. Dash {\em et al.}, A signal peptide missense mutation associated with nicotine dependence alters $\alpha$2*-nicotinic acetylcholine receptor function, {\em Neuropharmacology} {\bf 79}, 715--725 (2014).

\bibitem{BertrandEHTSBS2005}
D. Bertrand {\em et al.}, The CHRNB2 mutation I312M is associated with epilepsy and distinct memory deficits, {\em Neurobiology of Disease} {\bf 20} (3), 799--804 (2005).

\bibitem{VillaCMGMFCGBC2019}
C. Villa {\em et al.}, CHRNA2 and Nocturnal Frontal Lobe Epilepsy: Identification and Characterization of a Novel Loss of Function Mutation, {\em Frontiers in Molecular Neuroscience} {\bf 12} (2019).

\bibitem{LossiusSMBMSVCSPB2020}
K. Lossius {\em et al.}, Remarkable effect of transdermal nicotine in children with CHRNA4-related autosomal dominant sleep-related hypermotor epilepsy, {\em Epilepsy \& Behavior} {\bf 105}, 106944 (2020).

\bibitem{MazzaferroMCMS2022}
S. Mazzaferro {\em et al.}, Genetic Variant in Nicotinic Receptor $\alpha$4-Subunit Causes Sleep-Related Hyperkinetic Epilepsy via Increased Channel Opening, {\em International Journal of Molecular Sciences} {\bf 23} (20), 12124 (2022).

\bibitem{FuscoBPAGQBWC2000}
M. D. Fusco {\em et al.}, The nicotinic receptor $\beta$2 subunit is mutant in nocturnal frontal lobe epilepsy, {\em Nature Genetics} {\bf 26} (3), 275--276 (2000).

\bibitem{PhillipsFKZGHSSBBM2001}
H. A. Phillips {\em et al.}, CHRNB2 Is the Second Acetylcholine Receptor Subunit Associated with Autosomal Dominant Nocturnal Frontal Lobe Epilepsy*, {\em The American Journal of Human Genetics} {\bf 68} (1), 225--231 (2001).

\bibitem{DashL2014a}
B. Dash {\em et al.}, Two rare variations, D478N and D478E, that occur at the same amino acid residue in nicotinic acetylcholine receptor (nAChR) $\alpha$2 subunit influence nAChR function, {\em Neuropharmacology} {\bf 85}, 471--481 (2014).

\bibitem{YeungGSLL2025}
T. M. H. Yeung {\em et al.}, Slow-Channel Congenital Myasthenic Syndrome due to a CHRNA1 Variant Masquerading as Atypical Anorexia Nervosa: A Case Report, {\em Cureus} (2025).

\bibitem{GooneratneNMWCBGHSC2020}
I. K. Gooneratne {\em et al.}, Slow-Channel Congenital Myasthenic Syndrome due to a Novel Mutation in the Acetylcholine Receptor Alpha Subunit in a South Asian: A Case Report, {\em Journal of Neuromuscular Diseases} {\bf 8} (1), 163--167 (2020).

\bibitem{FinstererSSF2025}
J. Finsterer {\em et al.}, The new variant c.458C$>$G in CHRNA1 causes late-onset post-synaptic slow-channel congenital myasthenic syndrome, {\em Clinics} {\bf 80}, 100791 (2025).

\bibitem{GuanZHZZLWZL2025}
J. Guan {\em et al.}, Clinical and genetic features of congenital myasthenic syndrome due to the muscle acetylcholine receptor genes, {\em Brain and Development} {\bf 47} (5), 104412 (2025).

\bibitem{HodaGFPBAGRSMPBMSB2008}
J. Hoda {\em et al.}, Human Nocturnal Frontal Lobe Epilepsy: Pharmocogenomic Profiles of Pathogenic Nicotinic Acetylcholine Receptor $\beta$-Subunit Mutations outside the Ion Channel Pore, {\em Molecular Pharmacology} {\bf 74} (2), 379--391 (2008).

\bibitem{BoormanGS2000}
J. P. B. Boorman {\em et al.}, Stoichiometry of human recombinant neuronal nicotinic receptors containing the $\beta$3 subunit expressed in Xenopus oocytes, {\em The Journal of Physiology} {\bf 529} (3), 565--577 (2000).

\bibitem{Groot‐KormelinkBS2001}
P. J. Groot-Kormelink {\em et al.}, Formation of functional $\alpha$3$\beta$4$\alpha$5 human neuronal nicotinic receptors in Xenopus oocytes: a reporter mutation approach, {\em British Journal of Pharmacology} {\bf 134} (4), 789--796 (2001).

\bibitem{HallerDVMLASBJHSBG2011}
G. Haller {\em et al.}, Rare missense variants in CHRNB4 are associated with reduced risk of nicotine dependence, {\em Human Molecular Genetics} {\bf 21} (3), 647--655 (2011).

\bibitem{MannKHGSCNKMWKDCVITKEDSTBLBYMLEHKH2019}
N. Mann {\em et al.}, CAKUT and Autonomic Dysfunction Caused by Acetylcholine Receptor Mutations, {\em The American Journal of Human Genetics} {\bf 105} (6), 1286--1293 (2019).

\bibitem{McClureBegleyPSSLGXLP2014}
T. McClure-Begley {\em et al.}, Rare Human Nicotinic Acetylcholine Receptor $\alpha$4 Subunit (CHRNA4) Variants Affect Expression and Function of High-Affinity Nicotinic Acetylcholine Receptors, {\em The Journal of Pharmacology and Experimental Therapeutics} {\bf 348} (3), 410--420 (2014).

\bibitem{LenigerKHBBS2003}
T. Leniger {\em et al.}, A New Chrna4 Mutation with Low Penetrance in Nocturnal Frontal Lobe Epilepsy, {\em Epilepsia} {\bf 44} (7), 981--985 (2003).

\bibitem{SteinleinSMBSB2000}
O. K. Steinlein {\em et al.}, Independent Occurrence of the CHRNA4 Ser248Phe Mutation in a Norwegian Family with Nocturnal Frontal Lobe Epilepsy, {\em Epilepsia} {\bf 41} (5), 529--535 (2000).

\bibitem{JiangYYSSJ2018}
Y. Jiang {\em et al.}, CHRNA4 variant causes paroxysmal kinesigenic dyskinesia and genetic epilepsy with febrile seizures plus?, {\em Seizure} {\bf 56}, 88--91 (2018).

\bibitem{DibbensM2011}
L. Dibbens {\em et al.}, Genetic variations and associated pathophysiology in the management of epilepsy, {\em The Application of Clinical Genetics}, 113 (2011).

\bibitem{MazzaferroKNHS2024}
S. Mazzaferro {\em et al.}, Structural bases for stoichiometry-selective calcium potentiation of a neuronal nicotinic receptor, {\em British Journal of Pharmacology} {\bf 181} (13), 1973--1992 (2024).

\bibitem{TapiaKL2007}
L. Tapia {\em et al.}, Ca2+ Permeability of the ($\alpha$4)3($\beta$2)2 Stoichiometry Greatly Exceeds That of ($\alpha$4)2($\beta$2)3 Human Acetylcholine Receptors, {\em Molecular Pharmacology} {\bf 71} (3), 769--776 (2007).

\bibitem{JainKWKL2016}
A. Jain {\em et al.}, Unorthodox Acetylcholine Binding Sites Formed by $\alpha$5 and $\beta$3 Accessory Subunits in $\alpha$4$\beta$2* Nicotinic Acetylcholine Receptors, {\em Journal of Biological Chemistry} {\bf 291} (45), 23452--23463 (2016).

\bibitem{CadieuxDionMVTWBSGMCCBN2020}
M. Cadieux-Dion {\em et al.}, Variants inCHRNB2andCHRNA4Identified in Patients with Insular Epilepsy, {\em Canadian Journal of Neurological Sciences / Journal Canadien des Sciences Neurologiques} {\bf 47} (6), 800--809 (2020).

\bibitem{KuryatovBL2011}
A. Kuryatov {\em et al.}, Acetylcholine Receptor (AChR) $\alpha$5 Subunit Variant Associated with Risk for Nicotine Dependence and Lung Cancer Reduces ($\alpha$4$\beta$2)2$\alpha$5 AChR Function, {\em Molecular Pharmacology} {\bf 79} (1), 119--125 (2011).

\bibitem{DashCL2011}
B. Dash {\em et al.}, Reporter Mutation Studies Show That Nicotinic Acetylcholine Receptor (nAChR) $\alpha$5 Subunits and/or Variants Modulate Function of $\alpha$6*-nAChR, {\em Journal of Biological Chemistry} {\bf 286} (44), 37905--37918 (2011).

\bibitem{MartinelliCGLPPCDRBTCFCBCPBNKCGSTSRCLDTL2020}
S. Martinelli {\em et al.}, Co-occurring WARS2 and CHRNA6 mutations in a child with a severe form of infantile parkinsonism, {\em Parkinsonism \& Related Disorders} {\bf 72}, 75--79 (2020).

\bibitem{JensenHJ2013}
A. B. Jensen {\em et al.}, Elucidation of Molecular Impediments in the $\alpha$6 Subunit for in Vitro Expression of Functional $\alpha$6$\beta$4* Nicotinic Acetylcholine Receptors, {\em Journal of Biological Chemistry} {\bf 288} (47), 33708--33721 (2013).

\bibitem{ZouridakisGZCBT2014}
M. Zouridakis {\em et al.}, Crystal structures of free and antagonist-bound states of human $\alpha$9 nicotinic receptor extracellular domain, {\em Nature Structural \& Molecular Biology} {\bf 21} (11), 976--980 (2014).

\bibitem{ChaouchMGDSRLSWCNVMAAL2011}
A. Chaouch {\em et al.}, A retrospective clinical study of the treatment of slow-channel congenital myasthenic syndrome, {\em Journal of Neurology} {\bf 259} (3), 474--481 (2011).

\bibitem{OutteryckRLFZGEHVS2008}
O. Outteryck {\em et al.}, Novel   subunit mutation of the muscle acetylcholine receptor causing a slow-channel congenital myasthenic syndrome, {\em Journal of Neurology, Neurosurgery \& Psychiatry} {\bf 80} (4), 450--451 (2008).

\bibitem{ShenOMOTKGOE2016}
X. Shen {\em et al.}, Mutations Causing Slow-Channel Myasthenia Reveal That a Valine Ring in the Channel Pore of Muscle AChR is Optimized for Stabilizing Channel Gating: HUMAN MUTATION, {\em Human Mutation} {\bf 37} (10), 1051--1059 (2016).

\bibitem{CastroMGEE2007}
A. D. Castro {\em et al.}, Pathogenic point mutations in a transmembrane domain of the $\varepsilon$ subunit increase the Ca2+ permeability of the human endplate ACh receptor, {\em The Journal of Physiology} {\bf 579} (3), 671--677 (2007).

\bibitem{FucilePMME2000}
S. Fucile {\em et al.}, Human neuronal threonine-for-leucine-248 $\alpha$7 mutant nicotinic acetylcholine receptors are highly Ca 2+ permeable, {\em Proceedings of the National Academy of Sciences} {\bf 97} (7), 3643--3648 (2000).

\bibitem{NatarajanMCLBS2020}
K. Natarajan {\em et al.}, Mechanism of calcium potentiation of the $\alpha$7 nicotinic acetylcholine receptor, {\em Journal of General Physiology} {\bf 152} (9) (2020).

\bibitem{AndersenCSB2013}
N. Andersen {\em et al.}, Stoichiometry for activation of neuronal $\alpha$7 nicotinic receptors, {\em Proceedings of the National Academy of Sciences} {\bf 110} (51), 20819--20824 (2013).

\bibitem{LeeS2005}
W. Y. Lee {\em et al.}, Principal pathway coupling agonist binding to channel gating in nicotinic receptors, {\em Nature} {\bf 438} (7065), 243--247 (2005).

\bibitem{MihaylovaSGLKWSGHHAML2010}
V. Mihaylova {\em et al.}, Molecular characterisation of congenital myasthenic syndromes in Southern Brazil, {\em Journal of Neurology, Neurosurgery \& Psychiatry} {\bf 81} (9), 973--977 (2010).

\bibitem{KuryatovGNOL1997}
A. Kuryatov {\em et al.}, Mutation Causing Autosomal Dominant Nocturnal Frontal Lobe Epilepsy Alters Ca$^{2+}$Permeability, Conductance, and Gating of Human $\alpha$4$\beta$2 Nicotinic Acetylcholine Receptors, {\em The Journal of Neuroscience} {\bf 17} (23), 9035--9047 (1997).

\bibitem{LiangSMBDBD2004}
Y. Liang {\em et al.}, Functional polymorphisms in the human $\beta$4 subunit of nicotinic acetylcholine receptors, {\em Neurogenetics} {\bf 6} (1), 37--44 (2004).

\bibitem{MuellerHSJBRMTAL2006}
J. M{\"u}ller {\em et al.}, G.P.3 08 Novel mutations in the CHRNB1 gene in three patients affected by a congenital myasthenic syndrome, {\em Neuromuscular Disorders} {\bf 16} (9-10), 661 (2006).

\bibitem{ShenBESE2012}
X. Shen {\em et al.}, Highly fatal fast-channel syndrome caused by AChR $\varepsilon$ subunit mutation at the agonist binding site, {\em Neurology} {\bf 79} (5), 449--454 (2012).

\bibitem{SineSWOLTBBVE2002}
S. M. Sine {\em et al.}, Naturally Occurring Mutations at the Acetylcholine Receptor Binding Site Independently Alter ACh Binding and Channel Gating, {\em The Journal of General Physiology} {\bf 120} (4), 483--496 (2002).

\bibitem{EstephanSSTMOCSMCZAMHRATZ2018}
E. D. P. Estephan {\em et al.}, A common CHRNE mutation in Brazilian patients with congenital myasthenic syndrome, {\em Journal of Neurology} {\bf 265} (3), 708--713 (2018).

\bibitem{EirisPunalFGPLQBP2020}
J. M. Eir{\'i}s Pu{\~n}al {\em et al.}, Respuesta cl{\'i}nica y neurofisiol{\'o}gica a la efedrina en un paciente con s{\'i}ndrome miast{\'e}nico cong{\'e}nito de canal lento, {\em Revista de Neurolog{\'i}a} {\bf 71} (06), 221 (2020).

\bibitem{GomezMVNSCSRKVWL2001}
C. M. Gomez {\em et al.}, Novel delta subunit mutation in slow-channel syndrome causes severe weakness by novel mechanisms, {\em Annals of Neurology} {\bf 51} (1), 102--112 (2001).

\bibitem{ShenDSZNSSE2020}
X. Shen {\em et al.}, A novel fast-channel myasthenia caused by mutation in $\beta$ subunit of AChR reveals subunit-specific contribution of the intracellular M1-M2 linker to channel gating, {\em Experimental Neurology} {\bf 331}, 113375 (2020).

\bibitem{ShenMWBSSE2019}
X. Shen {\em et al.}, Slow-channel myasthenia due to novel mutation in M2 domain of AChR delta subunit, {\em Annals of Clinical and Translational Neurology} {\bf 6} (10), 2066--2078 (2019).

\bibitem{HallerLEHGS2014}
G. Haller {\em et al.}, Functional Characterization Improves Associations between Rare Non-Synonymous Variants in CHRNB4 and Smoking Behavior, {\em PLoS ONE} {\bf 9} (5), e96753 (2014).

\bibitem{SabatelliEACMLMLTSDGDMFLMODTNZ2009}
M. Sabatelli {\em et al.}, Rare missense variants of neuronal nicotinic acetylcholine receptor altering receptor function are associated with sporadic amyotrophic lateral sclerosis, {\em Human Molecular Genetics} {\bf 18} (20), 3997--4006 (2009).

\bibitem{ShenBSE2012}
X. Shen {\em et al.}, Myasthenic syndrome AChR$\alpha$ C-loop mutant disrupts initiation of channel gating, {\em Journal of Clinical Investigation} {\bf 122} (7), 2613--2621 (2012).

\bibitem{NooriMVSTMBMBFUSS2018}
H. R. Noori {\em et al.}, Alcohol reduces muscle fatigue through atomistic interactions with nicotinic receptors, {\em Communications Biology} {\bf 1} (1) (2018).

\bibitem{MukhtasimovaS2007}
N. Mukhtasimova {\em et al.}, An Intersubunit Trigger of Channel Gating in the Muscle Nicotinic Receptor, {\em The Journal of Neuroscience} {\bf 27} (15), 4110--4119 (2007).

\bibitem{MukhtasimovaFS2005}
N. Mukhtasimova {\em et al.}, Initial Coupling of Binding to Gating Mediated by Conserved Residues in the Muscle Nicotinic Receptor, {\em The Journal of General Physiology} {\bf 126} (1), 23--39 (2005).

\bibitem{LeeFS2008}
W. Y. Lee {\em et al.}, Nicotinic Receptor Interloop Proline Anchors $\beta$1-$\beta$2 and Cys loops in Coupling Agonist Binding to Channel Gating, {\em The Journal of General Physiology} {\bf 132} (2), 265--278 (2008).

\bibitem{ShenBNSE2016}
X. Shen {\em et al.}, Investigation of Congenital Myasthenia Reveals Functional Asymmetry of Invariant Acetylcholine Receptor (AChR) Cys-loop Aspartates, {\em Journal of Biological Chemistry} {\bf 291} (7), 3291--3301 (2016).

\bibitem{WebsterMSTBSB2012}
R. Webster {\em et al.}, A novel congenital myasthenic syndrome due to decreased acetylcholine receptor ion-channel conductance, {\em Brain} {\bf 135} (4), 1070--1080 (2012).

\bibitem{WebsterLCLB2014}
R. Webster {\em et al.}, Fast-channel congenital myasthenic syndrome with a novel acetylcholine receptor mutation at the $\alpha$--$\varepsilon$ subunit interface, {\em Neuromuscular Disorders} {\bf 24} (2), 143--147 (2014).

\bibitem{ShenOTBGSE2003}
X. Shen {\em et al.}, Mutation causing severe myasthenia reveals functional asymmetry of AChR signature cystine loops in agonist binding and gating, {\em Journal of Clinical Investigation} {\bf 111} (4), 497--505 (2003).

\bibitem{CunyKTYA2016}
H. Cuny {\em et al.}, Key Structural Determinants in the Agonist Binding Loops of Human $\beta$2 and $\beta$4 Nicotinic Acetylcholine Receptor Subunits Contribute to $\alpha$3$\beta$4 Subtype Selectivity of $\alpha$-Conotoxins, {\em Journal of Biological Chemistry} {\bf 291} (45), 23779--23792 (2016).

\bibitem{HoneMALLWPBA2015}
A. J. Hone {\em et al.}, $\alpha$-Conotoxins Identify the $\alpha$3$\beta$4* Subtype as the Predominant Nicotinic Acetylcholine Receptor Expressed in Human Adrenal Chromaffin Cells, {\em Molecular Pharmacology} {\bf 88} (5), 881--893 (2015).

\bibitem{PiccariDBGF2011}
V. Piccari {\em et al.}, Modulation of the Ca2+ permeability of human endplate acetylcholine receptor-channel, {\em Cell Calcium} {\bf 49} (4), 272--278 (2011).

\bibitem{ShenFOSE2008}
X. Shen {\em et al.}, Congenital myasthenia--related AChR $\delta$ subunit mutation interferes with intersubunit communication essential for channel gating, {\em Journal of Clinical Investigation} {\bf 118} (5), 1867--1876 (2008).

\bibitem{WangOMBEBMCES2000}
H. Wang {\em et al.}, Fundamental Gating Mechanism of Nicotinic Receptor Channel Revealed by Mutation Causing a Congenital Myasthenic Syndrome, {\em The Journal of General Physiology} {\bf 116} (3), 449--462 (2000).

\bibitem{OhnoWMBBNQPSE1996}
K. Ohno {\em et al.}, Congenital Myasthenic Syndrome Caused by Decreased Agonist Binding Affinity Due to a Mutation in the Acetylcholine Receptor $\varepsilon$ Subunit, {\em Neuron} {\bf 17} (1), 157--170 (1996).

\bibitem{NovielloGMCBBSH2021}
C. M. Noviello {\em et al.}, Structure and gating mechanism of the $\alpha$7 nicotinic acetylcholine receptor, {\em Cell} {\bf 184} (8), 2121--2134.e13 (2021).

\bibitem{BertrandGDBC1993}
D. Bertrand {\em et al.}, Mutations at two distinct sites within the channel domain M2 alter calcium permeability of neuronal alpha 7 nicotinic receptor., {\em Proceedings of the National Academy of Sciences} {\bf 90} (15), 6971--6975 (1993).

\bibitem{ZhangDZXXGHLC2015}
Q. Zhang {\em et al.}, Functional Impact of 14 Single Nucleotide Polymorphisms Causing Missense Mutations of Human $\alpha$7 Nicotinic Receptor, {\em PLOS ONE} {\bf 10} (9), e0137588 (2015).

\bibitem{CharpantierWHOHARFBF2005}
E. Charpantier {\em et al.}, $\alpha$7 Neuronal Nicotinic Acetylcholine Receptors Are Negatively Regulated by Tyrosine Phosphorylation and Src-Family Kinases, {\em The Journal of Neuroscience} {\bf 25} (43), 9836--9849 (2005).

\bibitem{TsunekiKTKTMMWKKS2007}
H. Tsuneki {\em et al.}, Novel G423S Mutation of Human $\alpha$7 Nicotinic Receptor Promotes Agonist-Induced Desensitization by a Protein Kinase C-Dependent Mechanism, {\em Molecular Pharmacology} {\bf 71} (3), 777--786 (2007).

\bibitem{BroideSJPPDD2002}
R. S. Broide {\em et al.}, Increased Sensitivity to Nicotine-Induced Seizures in Mice Expressing the L250T $\alpha$7 Nicotinic Acetylcholine Receptor Mutation, {\em Molecular Pharmacology} {\bf 61} (3), 695--705 (2002).

\bibitem{EllisonGWSMO2004}
M. Ellison {\em et al.}, $\alpha$-Conotoxins ImI and ImII Target Distinct Regions of the Human $\alpha$7 Nicotinic Acetylcholine Receptor and Distinguish Human Nicotinic Receptor Subtypes, {\em Biochemistry} {\bf 43} (51), 16019--16026 (2004).

\bibitem{GillThindDDSM2015}
J. K. Gill-Thind {\em et al.}, Structurally Similar Allosteric Modulators of $\alpha$7 Nicotinic Acetylcholine Receptors Exhibit Five Distinct Pharmacological Effects, {\em Journal of Biological Chemistry} {\bf 290} (6), 3552--3562 (2015).

\bibitem{HuangZSPZ2008}
X. Huang {\em et al.}, Modeling Binding Modes of $\alpha$7 Nicotinic Acetylcholine Receptor with Ligands: The Roles of Gln117 and Other Residues of the Receptor in Agonist Binding, {\em Journal of Medicinal Chemistry} {\bf 51} (20), 6293--6302 (2008).

\bibitem{Engel1996}
A. Engel, New mutations in acetylcholine receptor subunit genes reveal heterogeneity in the slow-channel congenital myasthenic syndrome, {\em Human Molecular Genetics} {\bf 5} (9), 1217--1227 (1996).

\bibitem{MichalkSBRPMBNJYOSWFWHWBMHH2008}
A. Michalk {\em et al.}, Acetylcholine Receptor Pathway Mutations Explain Various Fetal Akinesia Deformation Sequence Disorders, {\em The American Journal of Human Genetics} {\bf 82} (2), 464--476 (2008).

\bibitem{TanMX2016}
J. Tan {\em et al.}, A Missense Mutation in Epsilon-subunit of Acetylcholine Receptor Causing Autosomal Dominant Slow-channel Congenital Myasthenic Syndrome in a Chinese Family, {\em Chinese Medical Journal} {\bf 129} (21), 2596--2602 (2016).

\bibitem{AngeliniLSMV2018}
C. Angelini {\em et al.}, Clinical and genetic characterization of an Italian family with slow-channel syndrome, {\em Neurological Sciences} {\bf 40} (3), 503--507 (2018).

\bibitem{LorenzoniDAHFTLWKS2022}
P. J. Lorenzoni {\em et al.}, Congenital myasthenic syndrome in a cohort of patients with `double' seronegative myasthenia gravis, {\em Arquivos de Neuro-Psiquiatria} {\bf 80} (1), 69--74 (2022).

\bibitem{BoonJWKKVEV2022}
H. T. Boon {\em et al.}, Slow Channel Syndrome Revisited: 40 Years Clinical Follow-Up and Genetic Characterization of Two Cases, {\em Journal of Neuromuscular Diseases} {\bf 9} (4), 525--532 (2022).

\bibitem{WangMOSTBTBES1999}
H. Wang {\em et al.}, Acetylcholine receptor M3 domain: stereochemical and volume contributions to channel gating, {\em Nature Neuroscience} {\bf 2} (3), 226--233 (1999).

\bibitem{Shen2004}
X. Shen, Subunit-specific contribution to agonist binding and channel gating revealed by inherited mutation in muscle acetylcholine receptor M3-M4 linker, {\em Brain} {\bf 128} (2), 345--355 (2004).

\bibitem{ShenBSDPYVNDSE2018}
X. Shen {\em et al.}, Mutations causing congenital myasthenia reveal principal coupling pathway in the acetylcholine receptor $\varepsilon$-subunit, {\em JCI Insight} {\bf 3} (2) (2018).

\bibitem{ShenNMISOE2023}
X. Shen {\em et al.}, Impaired gating of $\gamma$- and $\varepsilon$-AChR respectively causes Escobar syndrome and fast-channel myasthenia, {\em Annals of Clinical and Translational Neurology} {\bf 10} (5), 732--743 (2023).

\bibitem{SineOBAMPE1995}
S. M. Sine {\em et al.}, Mutation of the acetylcholine receptor $\alpha$ subunit causes a slow-channel myasthenic syndrome by enhancing agonist binding affinity, {\em Neuron} {\bf 15} (1), 229--239 (1995).

\bibitem{MiloneWOPFSBGSE1998}
M. Milone {\em et al.}, Mode Switching Kinetics Produced by a Naturally Occurring Mutation in the Cytoplasmic Loop of the Human Acetylcholine Receptor $\varepsilon$ Subunit, {\em Neuron} {\bf 20} (3), 575--588 (1998).

\bibitem{AzumaNTSIIONAIOMNKSO2015}
Y. Azuma {\em et al.}, Congenital myasthenic syndrome in Japan: Ethnically unique mutations in muscle nicotinic acetylcholine receptor subunits, {\em Neuromuscular Disorders} {\bf 25} (1), 60--69 (2015).
